\documentclass[a4paper, twocolumn, twoside, 10pt]{article}
\usepackage[english]{babel}
\usepackage[
    inner=1cm,    % 内侧边距（靠近书脊）,稍微加大一点,防止字被装订遮住
    outer=1cm,  % 外侧边距（靠近手边）,保持原样
    top=1.5cm, 
    bottom=1.5cm,
]{geometry}
\usepackage{fontspec}
\usepackage{amsmath, amssymb} % 数学公式
\usepackage{ctex}             % 中文支持
\usepackage{xcolor}
\usepackage{listings}
\usepackage{microtype}
\usepackage[
    unicode,
    bookmarks=true,
    bookmarksopen=true
]{hyperref}
\setlength{\columnseprule}{0.4pt}
\setlength{\columnsep}{12pt}
% 代码块样式
\lstset{
    basicstyle=\ttfamily\small,
    backgroundcolor=\color{gray!10},
    frame=single,
    breaklines=true,
    xleftmargin=0.1\linewidth,
    xrightmargin=0.1\linewidth,
    breakatwhitespace=true,
    columns=flexible
}
\pagestyle{plain}

\title{\heiti 手动安装Linux} 

\begin{document}
% \maketitle
\section{Live环境基础设置}

\subsection{连接 Wi-Fi}

查看网络设备：
\begin{lstlisting}
ip link
\end{lstlisting}

输入以下命令进入 iwd：
\begin{lstlisting}
iwctl
\end{lstlisting}

扫描和获取无线网络：
\begin{lstlisting}
station wlan0 scan
station wlan0 get-networks
\end{lstlisting}

然后连接无线网络：
\begin{lstlisting}
station wlan0 connect <WiFi名称>
\end{lstlisting}

退出输入：
\begin{lstlisting}
exit
\end{lstlisting}

输入密码后输入以下命令查看网络是否连接成功：
\begin{lstlisting}
ip a
\end{lstlisting}


同步时间：
\begin{lstlisting}
timedatectl set-ntp true
\end{lstlisting}

配置镜像源：
\begin{lstlisting}
reflector -a 12 -c cn -f 10 -l 20 --sort rate(可能更换score) --save /etc/pacman.d/mirrorlist
\end{lstlisting}

更新数据库+安装密钥：
\begin{lstlisting}
pacman -Sy archlinux-keyring
\end{lstlisting}

\subsection{硬盘分区}

查看磁盘：
\begin{lstlisting}
lsblk -pf
\end{lstlisting}

不确定时，通过以下指令列出详细信息：
\begin{lstlisting}
fdisk -l /dev/sda(或其他磁盘)
\end{lstlisting}

还是不确定，通过以下命令查看空间：
\begin{lstlisting}
cfdisk /dev/sda(或其他磁盘)
\end{lstlisting}

第一次使用时会弹出选项选择，选择 GPT 分区表。左右键选择 NEW ，创建一个100MB的分区，再选择 TYPE ，把类型修改为 EFI System。如果列表没有，可能使用的是MBR分区，要修改为GPT。通过以下命令修改：
\begin{lstlisting}
fdisk /dev/sda
Command (m for help):g
Command (m for help):w
\end{lstlisting}

接着创建一个分区，使用剩余的空间，选择 NEW ，直接回车，类型为 Linux filesystem（不对时需要更改）。

分区结束后选择 Write，输入 yes 确认，Quit退出。

再次查看分区：
\begin{lstlisting}
lsblk -pf
\end{lstlisting}

格式化分区：
\begin{lstlisting}
mkfs.fat -F 32 /dev/sda1（efi）
mkfs.btrfs /dev/sda2（空闲文件）
\end{lstlisting}

\subsection{创建子卷}

挂载分区：
\begin{lstlisting}
mount -t btrfs /dev/sda2 /mnt
\end{lstlisting}

创建/和home子卷：
\begin{lstlisting}
btrfs subvolume create /mnt/@
btrfs subvolume create /mnt/@home
\end{lstlisting}

再次查看分区：
\begin{lstlisting}
lsblk -pf
\end{lstlisting}

现在是根分区挂载到了mnt，需要取消挂载：
\begin{lstlisting}
umount /mnt
\end{lstlisting}

挂载root子卷到/mnt：
\begin{lstlisting}
mount -t btrfs -o subvol=/@,compress=zstd /dev/sda2 /mnt
mount --mkdir -t btrfs -o subvol=/@home,compress=zstd /dev/sda2 /mnt/home
mount --mkdir /dev/sda1 /mnt/efi
\end{lstlisting}

cd到/mnt，查看是否有@和@home子卷。

\subsection{安装系统}

安装软件到指定的根目录下：
\begin{lstlisting}
pacstrap -K /mnt base base-devel linux(或linux-zen) linux-firmware btrfs-progs
\end{lstlisting}
\begin{lstlisting}
pacstrap /mnt networkmanager vim sudo intel-ucode(或amd-ucode)
\end{lstlisting}

生成fstab文件：
\begin{lstlisting}
genfstab -U /mnt > /mnt/etc/fstab
\end{lstlisting}

进入系统：
\begin{lstlisting}
arch-chroot /mnt
\end{lstlisting}

手动设置时区：
\begin{lstlisting}
ln -s /usr/share/zoneinfo/Asia/Shanghai /etc/localtime 
\end{lstlisting}

自动设置时区：
\begin{lstlisting}
timedatectl set-timezone Asia/Shanghai
\end{lstlisting}

查看时间是否正确：
\begin{lstlisting}
timedatectl
\end{lstlisting}

调整时间误差：
\begin{lstlisting}
hwclock --systohc
\end{lstlisting}

系统本地化设置：
\begin{lstlisting}
vim /etc/locale.gen 
"/"键搜索 en_US.UTF-8 和 zhCN.UTF-8 ，去掉前面的注释，:wq 保存退出。
\end{lstlisting}

生成本地化设置：
\begin{lstlisting}
locale-gen
\end{lstlisting}

编辑locale.conf文件，设置系统语言环境：
\begin{lstlisting}
vim /etc/locale.conf
插入:LANG=en_US.UTF-8
保存退出
\end{lstlisting}

设置主机名：
\begin{lstlisting}
vim /etc/hostname
插入:xxxxx（只能包含小写字母，数字0-9，短横杠-）
保存退出
\end{lstlisting}

修改root密码：
\begin{lstlisting}
passwd
输入完之后可以通过
cat /etc/shadow
查看密码是否修改成功，成功的话会显示加密后的密码。
\end{lstlisting}

安装引导程序：
\begin{lstlisting}
pacman -S grub efibootmgr
\end{lstlisting}

安装引导：
\begin{lstlisting}
grub-install --target=x86_64-efi --efi-directory=/efi --boot-directory=/efi --bootloader-id=GRUB(来一个你喜欢的名称)
\end{lstlisting}

如果主板没有识别，可以尝试以下命令：
\begin{lstlisting}
grub-install --target=x86_64-efi --efi-directory=/efi --boot-directory=/efi --removable
\end{lstlisting}

挂载链接：
\begin{lstlisting}
ln -s /efi/grub /boot/grub
\end{lstlisting}

生成grub配置文件：
\begin{lstlisting}
grub-mkconfig -o /efi(boot也可以)/grub/grub.cfg
\end{lstlisting}

配置双系统：
\begin{lstlisting}
pacman -S os-prober exfat-utils
\end{lstlisting}

编辑源文件：
\begin{lstlisting}
vim /etc/default/grub
shift + g 到文件末尾，插入以下内容：
GRUB_DISABLE_OS_PROBER=false
可选启动选项记忆功能：
取消注释GRUB_SAVEDEFAULT=true
gg到文件开头，修改：
GRUB_DEFAULT=saved
GRUB_CMDLINE_LINUX_DEFAULT="loglevel=5"
(显示日志)
保存退出
\end{lstlisting}

个人一个小问题，pice总线报错：
\begin{lstlisting}
GRUB_CMDLINE_LINUX_DEFAULT="loglevel=5 quiet splash pcie=noaer(或者pcie=nommconf)"
sudo update-grub
\end{lstlisting}

再次生成grub配置文件：
\begin{lstlisting}
grub-mkconfig -o /boot/grub/grub.cfg
\end{lstlisting}

安装Zram：
\begin{lstlisting}
pacman -S zram-generator
\end{lstlisting}

配置Zram：
\begin{lstlisting}
vim /etc/systemd/zram-generator.conf
插入以下内容：
[zram0]
zram-size = ram
compression-algorithm = zstd
保存退出
\end{lstlisting}

编辑源文件：
\begin{lstlisting}
vim /etc/default/grub
gg到文件开头，修改：
GRUB_CMDLINE_LINUX_DEFAULT="loglevel=5 quiet splash pcie=noaer zswap.enabled=0"
保存退出
\end{lstlisting}

再次生成grub配置文件：
\begin{lstlisting}
grub-mkconfig -o /boot/grub/grub.cfg
\end{lstlisting}

退出chroot环境：
\begin{lstlisting}
exit
\end{lstlisting}

重启：
\begin{lstlisting}
reboot
\end{lstlisting}

\section{系统初设置}

\subsection{网络设置}

开启NetworkManager服务：
\begin{lstlisting}
systemctl enable --now NetworkManager
\end{lstlisting}

打开网络连接：
\begin{lstlisting}
nmtui
移动到
Activate a connection
选择连接的网络
输入密码
完成连接。
\end{lstlisting}

来一点点装逼软件：
\begin{lstlisting}
pacman -S fastfetch cmatrix
\end{lstlisting}

\subsection{快照设置}

更新一波：
\begin{lstlisting}
pacman -Syu
\end{lstlisting}

来一些基本设置：
\begin{lstlisting}
vim /etc/environment
插入以下内容：
EDITOR=vim
保存退出
exit
重新登录让环境变量生效
\end{lstlisting}

创建普通用户：
\begin{lstlisting}
useradd -mG wheel username
passwd username
visudo
找到以下行，取消注释：
%wheel ALL=(ALL:ALL) ALL
保存退出
\end{lstlisting}

开启32位软件源：
\begin{lstlisting}
vim /etc/pacman.conf
找到以下行，取消注释：
[multilib] 
Include = /etc/pacman.d/mirrorlist
在最底部插入：
[archlinuxcn]
Server = https://mirrors.aliyun.com
（mirrors.ustc.edu.cn）/archlinuxcn/$arch
保存退出
pacman -Sy archlinuxcn-keyring
pacman -S yay paru
\end{lstlisting}

安装快照工具：
\begin{lstlisting}
pacman -S snapper snap-pac btrfs-assistant grub-btrfs inotify-tools
\end{lstlisting}

开启快照服务：
\begin{lstlisting}
systemctl enable --now grub-btrfsd
重启电脑
\end{lstlisting}

创建快照：
\begin{lstlisting}
snapper -c root create-config /
snapper -c home create-config /home
snapper -c root create --description "快照描述"
snapper -c home create --description "快照描述"
pacman -S linux-lts
grub-mkconfig -o /boot/grub/grub.cfg
\end{lstlisting}

恢复快照（root身份）：
\begin{lstlisting}
snapper -c root list
snapper -c root undochange (使用的快照序号)..0
或
btrfs-assistant -l
btrfs-assistant -r (使用的快照序号)
\end{lstlisting}

\subsection{驱动设置}

安装头文件：
\begin{lstlisting}
pacman -S --needed linux-zen-headers linux-lts-headers
\end{lstlisting}

安装显卡驱动：
\begin{lstlisting}
# 英伟达
pacman -S nvidia-open-dkms(去官网找对应的包) nvidia-utils lib32-nvidia-utils
# intel 
pacman -S mesa lib32-mesa vulkan-intel lib32-vulkan-intel
# amd
pacman -S mesa lib32-mesa vulkan-radeon lib32-vulkan-radeon xf86-video-amdgpu
\end{lstlisting}

视频编辑码的驱动：
\begin{lstlisting}
# intel
pacman -S intel-media-driver （或者libva-intel-driver）
# nvidia
pacman -S --needed nvidia-utils libva-nvidia-driver
# amd 不需要
reboot
\end{lstlisting}

登录普通账户
安装固件软件：
\begin{lstlisting}
sudo pacman -S sof-firmware alsa-ucm-conf alsa-firmware
sudo pacman -S pipewire wireplumber pipewire-alsa pipewire-pulse pipewire-jack
systemctl --user enable pipewire pipewire-pulse wireplumber
sudo pacman -S bulez
sudo systemctl enable --now bluetooth
sudo pacman -S power-profiles-daemon
sudo systemctl enable --now power-profiles-daemon
sudo pacman -S noto-fonts noto-fonts-emoji adobe-source-han-sans-cn-fonts
sudo pacman -S flatpak
sudo flatpak remote-modify flathub --url=https://mirrors.sjtu.edu.cn/flathub
reboot
\end{lstlisting}

受苦了，存档一下吧
\begin{lstlisting}
sudo snapper -c root create --description "初始设置完成"
sudo snapper -c home create --description "初始设置完成"
\end{lstlisting}

\section{桌面环境(niri篇)}

\subsection{安装桌面环境}

浅浅的安装一下：
\begin{lstlisting}
sudo pacman -S fish
fish
sudo pacman -S niri xwayland-satellite xdg-desktop-portal-gnome fuzzel kitty(alacritty) firefox wqy-zenhei
选择pipewire-jack作为默认音频后端:2
选择noto-fonts作为firefox默认字体:2
niri-session
\end{lstlisting}

super+shift+e退出
\begin{lstlisting}
sudo pacman -S --needed vim
vim ~/.config/niri/config.kdl
搜索alacritty，找到以下内容：
{spawn "alacritty";}
修改为kitty
想自动进入fish
{spawn "kitty" "-e" "fish";}
或
{spawn-sh "kitty -e fish";}
搜索output，找到以下内容：
/-output "eDP-1"
同时super+t再打开一个终端
niri msg outputs
找到接口名称(eDP-1)
找个分辨率+刷新率
回到ouput块，修改为：
output "eDP-1" {
    mode "1920x1080@60"
    scale 2.5
    postion x=0 y=0
    focus-at-startup
}
保存:w
\end{lstlisting}

设置取消鼠标加速：
\begin{lstlisting}
接着这个文件
搜索/mouse
mouse {
    accel-profile "flat"
    scroll-factor horizontal=2.0 vertical=-1.0
    自己设置试试，不需要可以进行注释
}
上方natural-scroll可以设置自然滚动
注释掉就是正常滚动。
保存退出
\end{lstlisting}

安装点重要程序：
\begin{lstlisting}
sudo pacman -S libnotify mako polkit-gnome
\end{lstlisting}

进入niri配置文件
\begin{lstlisting}
找到spawn-at-startup
直接写！
spawn-at-startup "waybar"
spawn-at-startup "/usr/lib/polkit-gnome
/polkit-gnome-authentication-agent-1"
spawn-at-startup "mako"
保存退出
mako & disown
/usr/lib/polkit-gnome
/polkit-gnome-authentication-agent-1 & disown
测试一下
notify-send hello
mkdir -p ~/.config/mako
vim ~/.config/mako/config
输入以下内容：
default-timeout=5000
border-radius=8
保存退出
makoctl reload
vim ~/.config/niri/config.kdl
gg来到顶部，写入：
environment {
    LC_MESSAGES "zh_CN.UTF-8"
}
保存退出
\end{lstlisting}

如果是自动化安装：
\begin{lstlisting}
sudo vim /etc/locale.gen
取消注释zh_CN.UTF-8
sudo locale-gen
super + shift + e退出
exit*2
登录
niri-session
\end{lstlisting}

想用中文运行某个软件:
\begin{lstlisting}
LANG=zh_CN.UTF-8 firefox
\end{lstlisting}

\subsection{输入法设置}

安装输入法：
\begin{lstlisting}
sudo pacman -S fcitx5-im fcitx5-rime rime-ice-pinyin-git
vim ~/.config/niri/config.kdl
gg来到顶部，写入：
environment {
    LC_MESSAGES "zn_CN.UTF-8"
    XMODIFIERS "@im=fcitx"
}
保存退出
mkdir -p ~/.local/share/fcitx5/rime
vim ~/.local/share/fcitx5/rime
/default.custom.yaml
插入：
patch:
   __include: rime_ice_suggestion:/
保存退出
fcitx5 -d
super+d打开 Fcitx 5 配置
搜索rime，添加rime到左边
附加组件 - classic
设置竖向显示
\end{lstlisting}

添加自启动：
\begin{lstlisting}
进入niri配置文件
插入：spawn-at-startup "fcitx5"
搜索binds
插入：
Mod+F1 { spawn-sh "pkill fcitx5 || fcitx5"; }
\end{lstlisting}

\subsection{桌面功能}

安装桌面功能：
\begin{lstlisting}
sudo pacman -S ffmpegthumbnailer gvfs-smb nautilus-open-any-terminal file-roller gnome-keyring gst-plugins-base gst-plugins-good gst-libav
sudo ln -s /usr/bin/kitty /usr/bin/gnome-terminal
vim ~/.config/kitty/kitty.conf
shell /usr/bin/fish
niri配置文件
插入：
Mod+E { spawn "nautilus"; }
\end{lstlisting}

niri的门户设置:
\begin{lstlisting}
mkdir -p ~/.config/xdg-desktop-portal/
vim ~/.config/xdg-desktop-portal/
xdg-desktop-portal.conf
插入：
[preferred]
default=gnome;gtk;
org.freedesktop.impl.portal.FileChooser=gtk
\end{lstlisting}

锁屏功能：
\begin{lstlisting}
sudo pacman -S swaylock-effects
mkdir -p ~/.config/swaylock
vim ~/.config/swaylock/config
插入：
screenshots
clock
indicator
indicator-radius=200
indicator-thickness=15
effect-blur=10x5

niri配置文件
插入：
Super+Alt+L hotkey-overlay-title="Lock the Screen: swaylock"{ spawn "swaylock"; }

sudo pacman -S swayidle
mkdir -p ~/.config/niri/scripts
vim ~/.config/niri/scripts/swayidle.sh
插入：
#!/usr/bin/env bash
swayidle -w \
    timeout 300 'swaylock -f ' \
    timeout 600 'niri msg action power-off-monitors' \
    resume 'niri msg action power-on-monitors' \
    timeout 1200 'systemctl suspend' 
保存退出
chmod +x ~/.config/niri/scripts/swayidle.sh

niri配置文件
插入：
spawn-at-startup "~/.config/niri/scripts/swayidle.sh"
保存退出

bash ~/.config/niri/scripts/swayidle.sh & disown
\end{lstlisting}

蓝牙功能：
\begin{lstlisting}
sudo pacman -S --needed blueberry bluez
sudo systemctl enable --now bluetooth
sudo pacman -S --needed power-profiles-daemon
sudo systemctl enable --now power-profiles-daemon
\end{lstlisting}


剪贴板：
\begin{lstlisting}
sudo pacman -S yay
yay -S wl-clipboard clipse-bin(/git) clipse-gui (报错就分开下)
clipse --listen
niri配置文件
插入：
spawn-sh-at-startup "clipse --listen"
Mod+Alt+V { spawn "clipse --gui"; }
保存一下

按上面的快捷键
命令行运行
niri msg pick-window
记住Title或者App ID：
Title: "Clipse GUI"
App ID: "clipse-gui"

niri配置文件
找到window-rule
window-rule {
    // This app-id regular expression will work for both:
    // - host Firefox (app-id is "firefox")
    // - Flatpak Firefox (app-id is "org.mozilla.firefox")
    match app-id=r#"firefox$"# title="^Picture-in-Picture$"
    match app-id="clipse-gui" (这个)
    open-floating true（和这个）
}
\end{lstlisting}

\subsection{屏幕功能}

截图功能：
\begin{lstlisting} 
niri配置文件
搜索screenshot
可以修改快捷方式
Mod+Alt+A { screenshot;}
保存位置在
screenshot-path "~/Pictures/Screenshots/Screenshot from %Y-%m-%d-%H-%M-%S.png"
保存退出
sudo pacman -S satty
mkdir -p ~/.config/satty
vim ~/.config/satty/config.toml
插入：
[general]
copy-command = "wl-copy"
focus-toggles-toolbars = true
actions-on-right-click = ["save-to-clipboard"]

cd /Templates
touch new

文件打开/.config/niri/scripts
右键一个文件
edit-screenshot.sh

vim编辑
#!/usr/bin/env bash

CLIPNOW=$(wl-paste | sha1sum)
niri msg action screenshot
while ["$wl-paste | sha1sum" = "$CLIPNOW"]; do
sleep .05
done
wl-paste | satty -f -

niri配置文件
Mod+shift+S { spawn "~/.config/niri/scripts/
edit-screenshot.sh"; }
\end{lstlisting}
3.5 
屏幕分享：
\begin{lstlisting}
sudo pacman -S --needed sof-firmware alsa-ucm-conf alsa-firmware
sudo pacman -S --needed pipewire wireplumber pipewire-alsa pipewire-pulse pipewire-jack
systemctl --user enable --now pipewire pipewire-pulse wireplumber
sudo pacman -S pavucontrol
niri配置文件
spawn-sh-at-startup "dbus-update-activation-environment --systemd WAYLAND_DISPLAY XDG_CURRENT_DESKTOP=niri & /usr/lib/xdg-desktop-portal-gnome"

sudo pacman -S --needed flatpak gnome-software
sudo flatpak remote-modify flathub --url=https://mirrors.sjtu.edu.cn/flathub

reboot
在软件商城安装obs
\end{lstlisting}

笔记本用户亮度切换：
\begin{lstlisting}
sudo pacman -S brightnessctl
\end{lstlisting}


\subsection{niri美化}

壁纸
\begin{lstlisting}
yay -S awww(老版是swww) waypaper
niri配置文件
spawn-at-startup "awww-daemon"
Mod+Alt+W {spawn "waypaper"}

match app-id="clipse-gui"下
加入 match app-id = "waypaper"开启浮动
\end{lstlisting}

任务栏
\begin{lstlisting}
sudo pacman -S waybar ttf-jetbrains-mono-nerd
niri配置文件
Mod+F2 {spawn-sh "pkill waybar || true && waybar";}
\end{lstlisting}

\section{任务栏diy}
\subsection{安装插件}

指令:
\begin{lstlisting}
sudo pacman -S wofi
\end{lstlisting}

\subsection{自定义}

自己创建waybar：
\begin{lstlisting}
mkdir -p ~/.config/waybar
touch ~/.config/waybar/config.jsonc
touch ~/.config/waybar/style.css
\end{lstlisting}

设置启动脚本：

这是我自己的waybar启动配置
\begin{lstlisting}
cd ~/.config/niri/scripts/
vim xxx.sh
#!/bin/bash
# 杀掉已有的 Waybar（避免重复启动）
# pkill waybar || true

# 启动左边 Waybar
waybar -c ~/.config/waybar/left_config.jsonc -s ~/.config/waybar/left_style.css &
sleep 0.3
# 启动右边 Waybar
waybar -c ~/.config/waybar/right_config.jsonc -s ~/.config/waybar/right_style.css &
\end{lstlisting}

\subsection{导入}





\end{document}